\documentclass[a4paper,12pt]{ctexart}
\usepackage{../../teaching-style}

\begin{document}
\pagestyle{empty}

\maintitle{多元微积分}
\begin{center}
  \large 扩展模块 · 偏导、重积分、矢量分析
\end{center}
\vspace{0.3cm}

\chaptertitle{第一章\quad 多元函数与偏导数}

\sectiontitle{1.1\quad 从一元到多元}

$z=f(x,y)$表示一个曲面。固定$x$看$y$，或固定$y$看$x$，
就回到了一元函数的情形。

\textbf{等高线}：$f(x,y)=C$——曲面上高度相同的点的集合，
就像地图上的等高线。

\sectiontitle{1.2\quad 偏导数}

$f_x(x_0,y_0)=\lim\limits_{\Delta x\to0}\dfrac{f(x_0+\Delta x,y_0)-f(x_0,y_0)}{\Delta x}$\\
固定$y$，对$x$求导。几何意义：沿着$x$方向切线的斜率。

类似地，$f_y(x_0,y_0)=\lim\limits_{\Delta y\to0}\dfrac{f(x_0,y_0+\Delta y)-f(x_0,y_0)}{\Delta y}$。

\begin{example}
$f(x,y)=x^2y+3xy^2$，求$f_x$和$f_y$。
\end{example}
\begin{solution}
$f_x=2xy+3y^2$，$f_y=x^2+6xy$。
\end{solution}

\sectiontitle{1.3\quad 全微分}

$dz = f_x\,dx + f_y\,dy$。
几何意义：用切平面近似曲面。
\[
\Delta z \approx f_x\Delta x + f_y\Delta y
\]

\sectiontitle{1.4\quad 高阶偏导}

$f_{xx}=\frac{\partial}{\partial x}(f_x)$，$f_{xy}=\frac{\partial}{\partial y}(f_x)$。
对于光滑函数，混合偏导与求导顺序无关：$f_{xy}=f_{yx}$。

\sectiontitle{1.5\quad 多元链式法则}

若$z=f(u,v)$，$u=u(x,y)$，$v=v(x,y)$，则：
\[
\frac{\partial z}{\partial x} = \frac{\partial z}{\partial u}\frac{\partial u}{\partial x} + \frac{\partial z}{\partial v}\frac{\partial v}{\partial x}
\]
树形图法可以帮助记忆。

\sectiontitle{习题1}

\begin{exercise}
\item $f(x,y)=3x^2y-2xy^3$，求$f_x$和$f_y$。
\item $f(x,y)=e^{xy}\ln y$，求$f_x$和$f_y$。
\item $f(x,y)=x^3y^2$，求$f_{xx}$、$f_{xy}$、$f_{yy}$。
\item $z=f(u,v)=u^2+v^2$，$u=x+y$，$v=x-y$，求$\frac{\partial z}{\partial x}$和$\frac{\partial z}{\partial y}$。
\end{exercise}

\newpage

\chaptertitle{第二章\quad 方向导数与梯度}

\sectiontitle{2.1\quad 为什么需要方向导数？}

偏导只告诉我们沿$x$或$y$方向的变化率。但如果我想知道
沿着任意方向$\vec u$的变化率呢？

方向导数：
\[
D_{\vec u}f = \lim_{t\to0}\frac{f(x_0+ta,y_0+tb)-f(x_0,y_0)}{t}
\]
其中$\vec u=(a,b)$为单位向量。

\sectiontitle{2.2\quad 方向导数的计算}

$D_{\vec u}f = f_x a + f_y b = \nabla f\cdot\vec u$。

这就引出了梯度。

\sectiontitle{2.3\quad 梯度——最陡上升方向}

\textbf{梯度}$\nabla f = (f_x,f_y)$是一个向量场。
它指向函数增长最快的方向，其模长等于最大变化率。

重要性质：梯度垂直于等高线。

\begin{center}
\begin{tikzpicture}[scale=0.6]
\draw[thick] (0,0) ellipse (2.5 and 1.5);
\draw[thick] (0,0) ellipse (2 and 1);
\draw[thick] (0,0) ellipse (1.5 and 0.5);
\draw[->,thick,red] (0.5,1.2) -- (1,1.8);
\draw[->,thick,red] (1.2,0.5) -- (2.2,0.8);
\draw[->,thick,red] (0,-1) -- (0,-2);
\node[red] at (2.5,2) {$\nabla f$};
\node at (1.5,-2.2){等高线$f(x,y)=C$};
\node[red] at (1.5,-1.5){梯度$\perp$等高线};
\end{tikzpicture}
\end{center}

\sectiontitle{2.4\quad 梯度下降法}

沿着负梯度方向走，函数值下降最快。
这是机器学习中优化算法的核心思想：
\[
x_{n+1} = x_n - \eta\nabla f(x_n)
\]

\sectiontitle{习题2}

\begin{exercise}
\item $f(x,y)=x^2+y^2$，求在$(1,2)$处沿方向$(1,1)$的方向导数。
\item 求$f(x,y)=x^2y$在$(1,2)$处的梯度。
\item 函数$f(x,y)=x^2+y^2$在点$(1,1)$处，沿哪个方向增长最快？最大变化率是多少？
\item 验证梯度垂直于等高线：$f(x,y)=x^2+y^2$在$(1,0)$处的梯度与过该点的等高线切线垂直。
\end{exercise}

\newpage

\chaptertitle{第三章\quad 重积分}

\sectiontitle{3.1\quad 二重积分}

$\iint_D f(x,y)\,dA$：把区域$D$切成小方块，每个方块用$f(x,y)$乘以面积，
再求和取极限。几何意义：曲面下的体积。

\sectiontitle{3.2\quad 直角坐标下的计算}

化为累次积分：
\[
\iint_D f(x,y)\,dA = \int_a^b\int_{g_1(x)}^{g_2(x)} f(x,y)\,dy\,dx
\]
或先对$x$再对$y$。关键是确定积分限——"从哪进、从哪出"。

\sectiontitle{3.3\quad 极坐标换元}

当区域是圆形或扇形时，极坐标更方便：
\[
x=r\cos\theta,\quad y=r\sin\theta,\quad dA = r\,dr\,d\theta
\]
\[
\iint_D f\,dA = \iint_{D'} f(r\cos\theta,r\sin\theta)\,r\,dr\,d\theta
\]

\sectiontitle{3.4\quad 三重积分}

$\iiint_V f\,dV$：类似二重积分，多了一维。

柱坐标：$x=r\cos\theta$，$y=r\sin\theta$，$z=z$，$dV=r\,dr\,d\theta\,dz$。
球坐标：$x=r\sin\phi\cos\theta$，$y=r\sin\phi\sin\theta$，$z=r\cos\phi$，
$dV=r^2\sin\phi\,dr\,d\phi\,d\theta$。

\sectiontitle{3.5\quad 应用}

求体积：$V=\iiint_V dV$。
求质量：$M=\iiint_V \rho\,dV$（$\rho$为密度）。
求质心：$\bar x = \frac1M\iiint_V x\rho\,dV$。
求转动惯量：$I=\iiint_V r^2\rho\,dV$。

\sectiontitle{习题3}

\begin{exercise}
\item 计算$\int_0^1\int_0^x (x+y)\,dy\,dx$。
\item 求$f(x,y)=x^2+y^2$在矩形$[0,1]\times[0,2]$上的二重积分。
\item 用极坐标计算$\iint_D (x^2+y^2)\,dA$，$D$为单位圆。
\item 求圆柱$x^2+y^2\le R^2$，$0\le z\le h$的体积。
\item 求半径为$R$的球体的体积（用球坐标）。
\end{exercise}

\newpage

\chaptertitle{第四章\quad 曲线积分与曲面积分}

\sectiontitle{4.1\quad 对弧长的曲线积分}

$\int_C f(x,y)\,ds$：沿曲线$C$累积$f$的值。
$ds=\sqrt{(dx)^2+(dy)^2}=\sqrt{1+(y')^2}\,dx$。

应用：求曲线形物体的质量、质心。

\sectiontitle{4.2\quad 对坐标的曲线积分}

$\int_C P\,dx+Q\,dy$：变力$\vec F=(P,Q)$沿曲线$C$所做的功。

与路径无关的条件：$\dfrac{\partial P}{\partial y}=\dfrac{\partial Q}{\partial x}$。

\sectiontitle{4.3\quad 对面积的曲面积分}

$\iint_S f(x,y,z)\,dS$：沿曲面累积。

\sectiontitle{4.4\quad 对坐标的曲面积分}

$\iint_S \vec F\cdot d\vec S$：向量场$\vec F$穿过曲面$S$的\textbf{通量}。
$d\vec S = \vec n\,dS$，$\vec n$为指向外侧的单位法向量。

\sectiontitle{习题4}

\begin{exercise}
\item 计算$\int_C (x^2+y^2)\,ds$，$C$是$y=x^2$从$(0,0)$到$(1,1)$的一段。
\item $\vec F=(y,x)$，沿$C$：$y=x^2$从$(0,0)$到$(1,1)$，求做功。
\item 判断$\int_C y\,dx+x\,dy$是否与路径无关，并计算从$(0,0)$到$(1,1)$沿任意路径的值。
\end{exercise}

\newpage

\chaptertitle{第五章\quad 梯度·散度·旋度}

\sectiontitle{5.1\quad 标量场与向量场}

\textbf{标量场}：空间中每个点对应一个数，如温度场$T(x,y,z)$。
\textbf{向量场}：空间中每个点对应一个向量，如流速场$\vec v(x,y,z)$。

\sectiontitle{5.2\quad 梯度 $\nabla f$}

梯度把标量场变成向量场：
\[
\nabla f = \left(\frac{\partial f}{\partial x},\;\frac{\partial f}{\partial y},\;\frac{\partial f}{\partial z}\right)
\]
指向函数增长最快的方向。

\sectiontitle{5.3\quad 散度 $\nabla\cdot\vec F$}

散度把向量场变成标量场：
\[
\nabla\cdot\vec F = \frac{\partial P}{\partial x} + \frac{\partial Q}{\partial y} + \frac{\partial R}{\partial z}
\]

\textbf{物理意义}：$\nabla\cdot\vec F>0$表示这个点是"源"（流体喷出），
$<0$表示是"汇"（流体吸入），$=0$表示不可压缩。

\sectiontitle{5.4\quad 旋度 $\nabla\times\vec F$}

旋度把向量场变成向量场：
\[
\nabla\times\vec F =
\begin{vmatrix}
\vec i & \vec j & \vec k\\
\frac{\partial}{\partial x} & \frac{\partial}{\partial y} & \frac{\partial}{\partial z}\\
P & Q & R
\end{vmatrix}
\]

\textbf{物理意义}：衡量向量场的旋转程度。$\nabla\times\vec F\neq\vec0$的地方有漩涡。

\begin{center}
\begin{tabular}{c|c|c}
算子 & 输入 & 输出\\ \hline
$\nabla f$（梯度） & 标量场 $f$ & 向量场 $\nabla f$\\
$\nabla\cdot\vec F$（散度） & 向量场 $\vec F$ & 标量场 $\nabla\cdot\vec F$\\
$\nabla\times\vec F$（旋度） & 向量场 $\vec F$ & 向量场 $\nabla\times\vec F$
\end{tabular}
\end{center}

\sectiontitle{习题5}

\begin{exercise}
\item $\vec F=(x^2y,\;y^2z,\;z^2x)$，求$\nabla\cdot\vec F$。
\item $\vec F=(-y,\;x,\;0)$，求$\nabla\times\vec F$。
\item 证明：$\nabla\times(\nabla f)=\vec0$（梯度场的旋度为0）。
\item 证明：$\nabla\cdot(\nabla\times\vec F)=0$（旋度场的散度为0）。
\item 流体流速$\vec v=(x,-y,0)$，判断哪些点是源、哪些是汇。
\end{exercise}

\newpage

\chaptertitle{第六章\quad 三大积分定理}

\sectiontitle{6.1\quad 格林定理（平面）}

\[
\oint_C P\,dx+Q\,dy = \iint_D \left(\frac{\partial Q}{\partial x} - \frac{\partial P}{\partial y}\right)\,dA
\]

左边是沿封闭曲线的环量，右边是内部旋度的累积。
\textbf{含义}：边界上的积分 = 内部的总变化。

特殊情形（取$P=-y$，$Q=x$）：$\oint_C x\,dy-y\,dx = 2\iint_D dA = 2S$。
可以用线积分求面积。

\sectiontitle{6.2\quad 高斯散度定理}

\[
\oiint_S \vec F\cdot d\vec S = \iiint_V \nabla\cdot\vec F\,dV
\]

左边是穿过封闭曲面的总通量，右边是内部所有散度的总和。
\textbf{含义}：曲面上的通量 = 内部'源'的总强度。

\sectiontitle{6.3\quad 斯托克斯定理}

\[
\oint_C \vec F\cdot d\vec r = \iint_S (\nabla\times\vec F)\cdot d\vec S
\]

左边是沿封闭曲线的环量，右边是曲面上旋度的通量。
\textbf{含义}：曲线上的环量 = 曲面上'漩涡'的总和。

\sectiontitle{6.4\quad 统一视角}

三个定理都是\textbf{微积分基本定理的推广}：
\[
\int_a^b f'(x)\,dx = f(b)-f(a)
\]
左边是"内部变化率的累积"，右边是"边界上的值"。
格林、高斯、斯托克斯都是这个思想在二维、三维的体现。

\begin{center}
\begin{tabular}{c|c|c}
定理 & 边界 & 内部\\ \hline
微积分基本定理 & 区间端点 $a,b$ & 区间内导数积分\\
格林定理 & 封闭曲线 $C$ & 内部区域旋度积分\\
高斯定理 & 封闭曲面 $S$ & 内部散度积分\\
斯托克斯定理 & 封闭曲线 $C$ & 曲面上旋度通量
\end{tabular}
\end{center}

\sectiontitle{6.5\quad 物理应用}

\textbf{麦克斯韦方程组（积分形式）}：
\begin{align*}
\oiint_S \vec E\cdot d\vec S &= \frac{Q}{\varepsilon_0}\quad\text{(高斯定律——电通量 = 内部电荷)} \\
\oiint_S \vec B\cdot d\vec S &= 0\quad\text{(磁高斯定律——无磁单极)} \\
\oint_C \vec E\cdot d\vec r &= -\frac{d\Phi_B}{dt}\quad\text{(法拉第定律——变化的磁场产生电场)} \\
\oint_C \vec B\cdot d\vec r &= \mu_0 I + \mu_0\varepsilon_0\frac{d\Phi_E}{dt}\quad\text{(安培-麦克斯韦定律)}
\end{align*}

每一个方程都是三大定理的直接应用。

\sectiontitle{习题6}

\begin{exercise}
\item 用格林定理计算$\oint_C y\,dx+x\,dy$，$C$为单位圆。
\item 验证高斯定理：$\vec F=(x,y,z)$，$S$为单位球面。
\item 用斯托克斯定理：$\vec F=(-y,x,0)$，$C$为单位圆。
\item 思考题：为什么$\nabla\times\vec E=0$的场（保守场）与路径无关？
\end{exercise}

\end{document}
